\newcommand{\tablereferenciados}{
 \autorescitacion. “\titulo”, \grado, \programa, Universidad de San Buenaventura \sede, \anio.
}

\renewcommand{\celdados}{ \multicolumn{2}{p{0.625\textwidth} }{\autorescitaciontexto. [1] }  }
\renewcommand{\celdacin}{  \multirow[t]{5}{0.68\textwidth}{  { \tablereferenciados}  }   }

\newcommand{\tablereferencias}{ Referencia/Reference }


\begin{tabular}
{       >{\centering}p{0.20\textwidth}     p{0.025\textwidth}    p{0.68\textwidth} }
\arrayrulecolor{orange}\toprule
Citar/How to cite         &                \celdados  \\ 
\cline{1-1}
\tablereferencias         & [1]          &\celdacin\\ 
                          &              &            \\
                          &              &            \\
Estilo/Style:             &              &            \\ 
IEEE (2014)               &              &            \\
\bottomrule
\end{tabular}


%\begin{tabular}{P{0.20\textwidth}p{0.74\textwidth}}
%  \noalign{\global\arrayrulewidth=0.4pt}
%  \arrayrulecolor{orange}\hline
%Citar/How to cite & \autorescitaciontexto. [1]\\
%\arrayrulecolor{orange}\hline
%
%Referencia/Reference\newline \newline
%Estilo/Style:\newline
%IEEE (2014)\newline
%&
%\begin{mylist} 
%\item[{[}1{]}] \autorescitacion. “Título” (Seleccione modalidad de grado nombre del programa académico), nombre de la facultad, Universidad de San Buenaventura Colombia, Ciudad, Año
%\end{mylist}\\ 
%
%\arrayrulecolor{orange}\hline
%\end{tabular}
